\documentclass[border=10pt]{standalone}
\usepackage{fontspec}
\setmainfont{Noto Serif}
\newfontfamily\bengalifont{Noto Serif Bengali}[Script=Bengali]
\usepackage[Latin,Bengali]{ucharclasses}
\setTransitionsFor{Bengali}{\begingroup\bengalifont}{\endgroup}
\usepackage{newunicodechar}
\newunicodechar{।}{{\bengalifont\char"0964}}
\newunicodechar{॥}{{\bengalifont\char"0965}}

\usepackage[table]{xcolor}
\usepackage{booktabs,tabularx,array}
\definecolor{bapilrule}{HTML}{1F2937}
\definecolor{bapilstripe}{HTML}{F3F4F6}
\renewcommand{\arraystretch}{1.28}
\setlength{\tabcolsep}{5pt}

\begin{document}
\begin{minipage}{17cm}
{\large\bfseries সারণি ৫।}\ {\large সক্ষমতার ঘাটতি বিশ্লেষণ}\par\vspace{7pt}
\small
\rowcolors{2}{bapilstripe}{white}
\begin{tabularx}{\linewidth}{@{}>{\hsize=0.876\hsize\raggedright\arraybackslash}X>{\hsize=1.237\hsize\raggedright\arraybackslash}X>{\hsize=1.237\hsize\raggedright\arraybackslash}X>{\hsize=1.237\hsize\raggedright\arraybackslash}X>{\hsize=0.412\hsize\raggedright\arraybackslash}X@{}}
\arrayrulecolor{bapilrule}\toprule
\rowcolor{white}\textbf{সক্ষমতা} & \textbf{বর্তমান অবস্থা} & \textbf{ঘাটতি} & \textbf{ব্লুপ্রিন্টের সাড়া} & \textbf{অধ্যায়} \\
\midrule
সরকারি নথিতে প্রবেশাধিকার & বহু সাইটে প্রকাশিত, প্রায়ই স্ক্যানকৃত PDF & বাংলা ও ইংরেজিতে একীভূত, অনুসন্ধানযোগ্য, উদ্ধৃতিযোগ্য সংগ্রহ নেই & নিয়ন্ত্রিত P0 কর্পাস ও নীতি অনুসন্ধান (P1) & ৮, ১২ \\
বাংলা ভাষা প্রযুক্তি & গবেষণা-মডেল ও ডেটাসেট বিদ্যমান [46, 48, 49, 81] & নীতি-ক্ষেত্রের বেঞ্চমার্ক নেই; OCR ও প্রমিতকরণ অসম & BanglaPolicyQA, প্রমিতকরণ নিয়ম, পর্যায়ভিত্তিক BPFM & ৯ \\
এআই শাসনের চর্চা & খসড়া নীতির শ্রেণি; নতুন উপাত্ত-আইন & ঝুঁকি-স্তর, গেট বা স্বচ্ছতা নিবন্ধনের কার্যকর ব্যবহার নেই & ঝুঁকি-স্তর, গেট, ERRC, নিবন্ধন & ১৩ \\
ডেটা তত্ত্বাবধান & উন্মুক্ত ডেটা পোর্টাল ও পরিসংখ্যান প্রকাশনা [31, 33] & মেটাডেটা, লাইসেন্স ও মানে অসামঞ্জস্য [78, 79] & DCAT মেটাডেটা, মানের মাপকাঠি, NRDEX চুক্তি & ৮ \\
কম্পিউটিং & বিশ্ববিদ্যালয়ে দুষ্প্রাপ্য [80] & জনস্বার্থমূলক এআই-এর জন্য যৌথ ও নিয়ন্ত্রিত পরিবেশ নেই & পরিমিত পাইলট হোস্টিং; অভিযোজনে যৌথ জাতীয় কম্পিউটিং & ১০ \\
সরকারে দক্ষতা & ডিজিটাল সেবায় শক্তিশালী অভিজ্ঞতা & এআই সাক্ষরতা ও নিশ্চয়তা-দক্ষতা সীমিত & ছয়টি প্রশিক্ষণ কর্মসূচি ও ফেলোশিপ & ৬ \\
মূল্যায়ন-সংস্কৃতি & উৎপাদ-কেন্দ্রিক প্রকল্প প্রতিবেদন & এআই সরঞ্জামের প্রকাশিত ও পুনরুৎপাদনযোগ্য মূল্যায়ন কম & বেঞ্চমার্ক, KPI, স্বাধীন মূল্যায়ন, DOI & ১৮ \\
সমন্বয় & একাধিক জাতীয় ডিজিটাল ও এআই কর্মসূচি & দ্বৈততা ও অস্পষ্ট ভূমিকার ঝুঁকি & কর্মবিভাজনের সমঝোতা স্মারক ও অংশীজন পরিকল্পনা & ৩, ৫ \\
\bottomrule
\end{tabularx}
\par\vspace{3pt}{\scriptsize\color{gray} উৎস: BAPIL কারিগরি ব্লুপ্রিন্ট, সংস্করণ ১.০ (২০২৬)। CC BY 4.0।}
\end{minipage}
\end{document}
